\documentclass[11pt,letterpaper]{article}
\usepackage[T1]{fontenc}
\usepackage[utf8]{inputenc}
\usepackage{lmodern,microtype}
\usepackage[margin=0.79in,headheight=15pt]{geometry}
\usepackage{parskip,enumitem,longtable,tabularx,booktabs,array,xcolor,fancyhdr,hyperref}
\definecolor{navy}{HTML}{17314C}
\definecolor{teal}{HTML}{087F78}
\hypersetup{colorlinks=true,linkcolor=navy,urlcolor=teal}
\setlength{\emergencystretch}{3em}
\setlist[itemize]{leftmargin=*,nosep,topsep=5pt}
\setlist[enumerate]{leftmargin=*,itemsep=3pt}
\renewcommand{\arraystretch}{1.22}
\pagestyle{fancy}\fancyhf{}
\fancyhead[L]{\small Enterprise Data Modeling}\fancyhead[R]{\small Documentation Suite}
\fancyfoot[C]{\thepage}
\newcommand{\srcbook}{\textit{The Data Model Resource Book}, revised edition, Volume 2}
\newcommand{\ResponseLine}{\par\vspace{1.05em}\noindent\rule{\linewidth}{0.4pt}\par}
\newcommand{\ResponseArea}{\ResponseLine\ResponseLine\ResponseLine}
\begin{document}
\begin{center}
{\LARGE\bfseries\color{navy} Logical to Physical Schema Design}\par
\vspace{0.25em}{\large Guided workbook | conceptual and operational analysis}\par
\vspace{0.3em}{\small Original engineering study material \textbullet{} 2026 edition}
\end{center}
\vspace{0.8em}
\section{Workbook objective}
Complete independently before looking at any model answer. Record assumptions, counterexamples and a proof of each chosen invariant. This is supplementary material; no exercise copies a worked example from the source book.


\section{Problem A: cardinality}
Given: A Party has zero or many AgreementParticipation rows; an Agreement has at least one participant. Participation may change role over time.

Tasks: Give a normalized SQL structure and explain which mandatory participation rule is difficult to express with simple foreign keys.

\subsection{Your working}
\ResponseArea\ResponseArea
\subsection{Verification notes}
\ResponseArea

\section{Problem B: alternate keys}
Given: An external product code can be reused by different sellers but not twice by the same seller on an overlapping effective period.

Tasks: Specify uniqueness scope, period semantics, constraint mechanisms and a concurrent-write test.

\subsection{Your working}
\ResponseArea\ResponseArea
\subsection{Verification notes}
\ResponseArea

\section{Problem C: update anomalies}
Given: A single flat record contains three addresses and three contact channels for each party.

Tasks: Design repeatable child entities, keys, cardinalities and a null-handling policy; justify whether the original record can be reconstructed.

\subsection{Your working}
\ResponseArea\ResponseArea
\subsection{Verification notes}
\ResponseArea

\section{Self-assessment}
\begin{itemize}
\item Can another reviewer reproduce the decomposition or integrity argument?
\item Did the model document what would falsify its assumptions?
\item Can every proposed invariant be linked to executable or review evidence?
\end{itemize}

\vfill
\hrule
{\footnotesize\textbf{Provenance.} Supplementary engineering practice: original content; the cited book provides domain-modeling context.}
\end{document}
